\documentclass[11pt]{article}
\usepackage[a4paper,margin=25mm]{geometry}
\usepackage{amsmath,amssymb,amsthm,booktabs,hyperref,xurl}
\setlength{\emergencystretch}{2em}
\hypersetup{hidelinks}
\newtheorem{theorem}{Theorem}
\newtheorem{proposition}[theorem]{Proposition}
\newtheorem{corollary}[theorem]{Corollary}
\newcommand{\Q}{\mathbb Q}\newcommand{\Z}{\mathbb Z}
\title{Three exact boundaries for further witness compression}
\author{\(\Z\Q\) Diophantine Project}
\date{23 November 2026}
\begin{document}\maketitle
\begin{abstract}
We isolate three precise limitations of natural compression attempts.
The existing 139-coordinate circuit for the fixed power $b^{5^{60}}$
is optimal in a triangular quadratic binomial model even when negative
and fractional Laurent exponents are allowed. A single quadratic
integer-valued test cannot recognize an unrestricted rational integer
tuple of length at least two. Finally, an explicit Pell family shows
that summing the private ordinary-$abc$ height expressions into one
common expression destroys the height contraction needed by the old
soundness argument. These are proved boundaries for specified methods,
not lower bounds for arbitrary Diophantine witnesses. No variable-count
improvement or new undecidability assertion is claimed by this note.
\end{abstract}

\section{A Laurent-exponent lower bound for the fixed-power circuit}

The initial available values are $1,b,b^2$, with $b>0$ allowed to vary
in an open interval. A \emph{triangular quadratic binomial step} adjoins
one nonzero coordinate $v$ and one relation between two monomials of
total degree at most two in $v$ and the preceding coordinates. Both
monomial coefficients are nonzero constants independent of $b$.
Every coordinate is required to have the form $c b^e$, with $c\ne0$
constant and $e\in\Q$. A relation from which the new coordinate
cancels entirely is not a determining step. Negative exponents,
fractional exponents, repeated old factors, and constant monomials are
all permitted. The model includes the ceiling-binary power circuit
of~\cite{ceiling}, but not arbitrary auxiliary rational functions or
simultaneously implicit blocks.

\begin{theorem}\label{width}
If the old exponent interval is $[m,M]$, with $m\le0\le M$, one
quadratic binomial step enlarges its width by at most a factor two.
Starting from $\{0,1,2\}$, after $L$ new coordinates the interval
width is therefore at most $2^{L+1}$.
\end{theorem}
\begin{proof}
All values are nonzero, so common monomial factors may be cancelled.
After cancellation the new coordinate occurs on just one side. Its
remaining multiplicity is either one or two. An identity between
$c_1b^{e_1}$ and $c_2b^{e_2}$ on an open interval forces
$e_1=e_2$. Fixed nonzero coefficients do not change this conclusion.

If the multiplicity is two, the other side contains at most two old
factors. Missing factors have exponent zero. The new exponent is
$(e_i+e_j)/2$, which belongs to $[m,M]$.

If the multiplicity is one, its side contains at most one old factor,
and the other side at most two. Again use zero exponents for missing
factors. The new exponent is
\[
                        e=e_i+e_j-e_k.
\]
Thus $2m-M\le e\le2M-m$. If $e\in[m,M]$, the interval is
unchanged. If $e>M$, its new width is
$e-m\le2(M-m)$; if $e<m$, its new width is
$M-e\le2(M-m)$. These cases prove the assertion. They include
repeated factors and all signs of the old exponents. A cancelled
relation with no new coordinate either supplies an old consistency
condition or is inconsistent; it cannot define a new exponent.

The initial width is two. Iterating the bound gives $2^{L+1}$.
A zero coefficient, a zero coordinate, or a coefficient depending on
$b$ would fall outside the stated model; none is silently used here.
\end{proof}

\begin{corollary}\label{optimal}
In this model, computing $b^{5^{60}}$ with $b^2$ initially available
requires at least 139 new coordinates. The existing circuit attains
this bound, so 139 is optimal in this model.
\end{corollary}
\begin{proof}
The exact inequalities are
\[
 2^{139}<5^{60}=867361737988403547205962240695953369140625<2^{140}.
\]
The exponent zero remains available, so an interval containing the
positive target exponent has width at least $5^{60}$. With $L\le138$,
Theorem~\ref{width} gives width at most $2^{139}$, a contradiction.

For the upper bound, adjoin $b^3$ and $b^5$ by
$b^3=b\,b^2$ and $b^5=b^3b^2$. Put
\[
 e_j=\left\lceil\frac{5^{60}}{2^{137-j}}\right\rceil
 \quad(0\le j\le137).
\]
Then $e_0=5$, $e_{137}=5^{60}$, and
$e_j=2e_{j-1}-\varepsilon_j$ with
$\varepsilon_j\in\{0,1\}$. Introduce the 137 subsequent powers by
\[
 v_j=v_{j-1}^2\quad(\varepsilon_j=0),\qquad
 b v_j=v_{j-1}^2\quad(\varepsilon_j=1).
\]
There are $2+137=139$ new coordinates. The guard $b\ne0$ gives a
unique rational lift in either case. The earlier complete system
supplies that guard as proved in~\cite{ceiling}.
\end{proof}

Both $b^3$ and $b^5$ are counted here; only $1,b,b^2$ are free
initial data. If additional powers are supplied without charging their
construction, the lower-bound input must be changed accordingly.
This corollary does not exclude a more economical construction using
nonmonomial intermediate values, a nontriangular system, the remaining
source equations, or a different universal source.

\section{One quadratic integer test cannot replace a whole tuple}

\begin{theorem}\label{quad}
Let $r\ge2$, and let $f\in\Q[X_1,\ldots,X_r]$ have total degree
at most two. Suppose $f(\Z^r)\subseteq\Z$. Then there is
$x\in\Q^r\setminus\Z^r$ with $f(x)\in\Z$.
Thus $f(x)\in\Z$ cannot, by itself, characterize $x\in\Z^r$.
\end{theorem}
\begin{proof}
If $f$ is constant, the assertion is immediate. If its degree is one,
choose $a\in\Z^r$ and a nonzero rational vector in the kernel of its
linear part. The resulting rational line through $a$ has constant
value $f(a)\in\Z$ and contains nonintegral rational points.

Suppose $f$ has degree two. Choose $a\in\Z^r$ with
$\nabla f(a)\ne0$. Such a point exists because a nonzero rational
polynomial cannot vanish on the whole integer grid. Translation gives
\[
                   f(a+x)-f(a)=L(x)+Q(x),
\]
where $L$ is a nonzero rational linear form and $Q$ a nonzero
rational homogeneous quadratic form. For every rational direction $v$
with $L(v)Q(v)\ne0$, the point
\[
                   x(v)=-\frac{L(v)}{Q(v)}v
\]
satisfies $L(x(v))+Q(x(v))=0$.

Choose a rational direction $v_0$ for which this denominator and
numerator are nonzero. On a sufficiently small closed real neighborhood
of $v_0$ both remain bounded away from zero, and $x(v)$ is bounded.
There are infinitely many distinct rational projective directions in
this neighborhood because $r\ge2$. Their images $x(v)$ are nonzero
and have those same distinct directions, so they are distinct.
Thus the level set $f=f(a)$ contains infinitely many bounded rational
points $a+x(v)$. Only finitely many of these can be integer points.
A nonintegral one has $f$-value $f(a)\in\Z$, as required.
\end{proof}

For the usual sum of squares, the familiar exact example is
$(3/5)^2+(4/5)^2=1$. The theorem covers every rational quadratic
polynomial that is integer-valued on the integer grid, including affine
and degenerate forms. It concerns an unrestricted tuple and one test.
It says nothing about higher-degree tests, several tests, quantified
auxiliary coordinates, or the special subvariety cut out by the outer
universal system.

\section{A shared height sum loses the scalar contraction}

The ordinary-$abc$ integer predicate uses the private expression
\[
                    M_i=B+2(X_i^8+T_i^2).
\]
Its soundness proof applies a two-term mixed-power height inequality;
see~\cite{abc56,signed}. A tempting coordinate saving would replace all
private expressions by
\begin{equation}\label{commonM}
                    M=B+2\sum_{i=1}^k(X_i^8+T_i^2).
\end{equation}
The following elementary obstruction is independent of $abc$.
As usual, $h(A/D)=\log\max(|A|,D)$ for a reduced rational number
with positive denominator.

\begin{proposition}\label{height}
Fix $k\ge3$ and $B\in\Q$. There are rational tuples with every
$X_i>0$, with $M=B+2k$ in~\eqref{commonM}, and with
$h(X_1)$ arbitrarily large. In particular no inequality
\[
                  \max_i h(X_i)\le C h(M)+C'
\]
with fixed finite constants $C,C'$ holds for all such rational tuples.
\end{proposition}
\begin{proof}
Start with $(n_0,t_0)=(1,0)$ and recursively set
\[
 (n_{j+1},t_{j+1})=(3n_j+4t_j,\;2n_j+3t_j).
\]
Direct expansion preserves $n_j^2-2t_j^2=1$. For $j\ge1$ both
coordinates are positive, and $n_j\ge3^j$, so $n_j$ is unbounded.
Write $n=n_j$, $t=t_j$ and set
\[
 X_1=1/n,\qquad X_2=\cdots=X_k=1,
\]
\[
 T_1=\frac{n^4-1}{n^4},\quad
 T_2=\frac{2nt}{n^4},\quad
 T_3=\frac{2t}{n^4},\qquad T_i=0\quad(i>3).
\]
The norm identity gives the exact sum of three squares
\[
 (n^4-1)^2+(2nt)^2+(2t)^2
 =n^8-2n^4+1+2(n^2-1)(n^2+1)=n^8-1.
\]
Consequently $\sum_iT_i^2=1-n^{-8}$ and
$\sum_iX_i^8=k-1+n^{-8}$. Their sum is exactly $k$, whereas
$h(X_1)=\log n\longrightarrow\infty$.
\end{proof}

The first three $T_i$ in this family are strictly positive. The
counterexample invalidates an unrestricted scalar height bound for
\eqref{commonM}. It does \emph{not} assert that these $T_i$ occur as
elliptic traces of valid point charts, or that the whole modified
Diophantine system has a nonintegral solution. A sharing theorem
exploiting those extra elliptic or source constraints could still exist;
it would require a new proof rather than the former two-term bound.

\section{Exact reproducibility and status}

The adjacent \texttt{verify\_compression\_barriers.py} independently
reconstructs the 139-step exponent sequence, checks the strict integer
inequalities in Corollary~\ref{optimal}, exercises the interval lemma
on signed rational edge cases, checks the quadratic level-set identity
symbolically, and verifies the first twelve Pell height counterexamples
by exact rational arithmetic. Its deterministic JSON output is recorded
in \texttt{certificate.json}. Finite examples do not replace the
universal proofs above. No complete coefficient table or result census
is modified, and no arithmetic premise is added.

\begin{thebibliography}{9}
\bibitem{ceiling} \(\Z\Q\) Diophantine Project,
\emph{A quadratic ceiling-binary circuit for a fixed power},
\nolinkurl{undecidability/papers/ceiling_power_quartic_checkpoint.tex}.
\bibitem{abc56} \(\Z\Q\) Diophantine Project,
\emph{Paired contact coordinates: an ordinary-abc quartic in 56 variables},
\nolinkurl{papers/quartic_existence/paired_contact56_20261119_v81.tex}.
\bibitem{signed} \(\Z\Q\) Diophantine Project,
\emph{Signed mixed-power heights and single-cover criteria},
\nolinkurl{papers/elliptic_heights/signed_mixed_power_height_20261117_v58.tex}.
\end{thebibliography}
\end{document}
